\section{The first wild index at five by residual comparison}\label{sec:exact-five}
This section gives a proof of the first wild index at five by separate residual comparisons; Section~\ref{sec:first-wild-general} proves the general-prime theorem from the graph valuation bound.
\begin{theorem}\label{thm:exact-five}
For the height-two filtration class over complete fields with perfect residue field, $e_K=20$, and $\mu_5\subset K$, one has
\[
                         \nuMin(5,20,2)=2.
\]
Thus $125$-torsion determines the finite special character uniformly, and $25$-torsion does not.
\end{theorem}
\begin{proof}
For the lower bound apply Theorem~\ref{thm:general-arc} with $p=5$, $n=2$, $e=20$, $A=1$. Its weights are $d=480$ and $b=380$, and its tangent multiplier is a primitive 24th root. The construction gives a supersingular group and a group filtered by height-one groups with the same 25-torsion but different special characters. This is an actual example over a base of the specified index, so the uniform index is at least two.

For the upper bound, take a comparison of 125-torsion and pass to the completed maximal unramified base. The determinant on inertia has the form $\chi_5^d\epsilon$, with $0\le d\le2$ and finite $\epsilon\in\mu_4$. Its wild cyclotomic order modulo 125 is at least five, because $v_5(e_K)=1$. Equality of the determinants therefore gives the same $d$ and the same finite factor. As in the proof of Theorem~\ref{thm:second-wild-exact}, dimensions zero and two are determined residually, and the dimension-one case consists of filtered height-one groups or whole potentially supersingular groups. Two filtered groups are determined residually. A mixed comparison would require $e_K>5^{2\cdot3-3-1}=25$ by Proposition~\ref{prop:mixed-obstruction}, contrary to $e_K=20$.

For two supersingular groups, their common residual inertia is either irreducible, tame reducible, or wild reducible. Theorems~\ref{thm:irreducible-recovery}, \ref{thm:tame-five-recovery}, and \ref{thm:wild-five-recovery} settle these cases. Passing to the completed unramified base does not change the inertia character being compared, so the conclusion holds for the original perfect residue field. This proves the upper bound two.
\end{proof}
